\begin{proof}[Proof of \cref{obj:lem:factorial-moment-mean}]
Write \(M=M(n)\), \(B=B(n)\), and let \(q=q_k\) be the four-cell vector for category \(k\). Set
\[
g_{M,j}=(-1)^j\frac{2^{2j+3}}{M(M+j+2)}\binom{M+j+2}{2j+4}.
\]
The polynomial whose mean is evaluated is the finite polynomial
\[
P_{M,B}(u)
=
B^{-1}s(u)t_1(u)\sum_{j=0}^{M-2}g_{M,j}\{B^{-1}s_1(u)\}^j
-
B^{-1}s(u)t_0(u)\sum_{j=0}^{M-2}g_{M,j}\{B^{-1}s_0(u)\}^j .
\]
We compute the factorial lift in its sparse arm expansion.  The factorial-moment contribution is
\[
\widehat P_k=\widehat A_{1,k}-\widehat A_{0,k},
\]
where, for \(a\in\{0,1\}\),
\[
\widehat A_{a,k}
=
\sum_{j=0}^{M-2}\sum_{t=0}^{j}\sum_{c\in\{0,1\}^2}
B^{-1}g_{M,j}(B^{-1})^j\binom jt\,
\frac{\prod_{b,y}(N^{(1)}_{bky})_{r^{a,c,j,t}_{by}}}{(m_1)_{|r^{a,c,j,t}|}}.
\]
The merged exponent vector \(r^{a,c,j,t}\) is defined by
\[
u^{r^{a,c,j,t}}
=
u_c\,u_{a1}\,u_{a0}^{\,t}\,u_{a1}^{\,j-t}.
\]
Thus \(\widehat P_k\) is the difference of the two displayed sparse arm sums.  % lean: factorialPolynomialContribution_eq_sparseArms

The cutoff condition \(n\ge N_0\) includes \(4M^2\le m_1\); since \(M>0\), this implies \(M\le m_1\). In every displayed summand, \(0\le t\le j\le M-2\), and the merged exponent vector has total degree
\[
|r^{a,c,j,t}|=j+2\le M\le m_1 .
\]
For each such multi-index \(r\), the factorial-moment identity for the second-split multinomial cell counts gives
\[
\mathbb E_{P^{\otimes\infty}}\!\left[
\frac{\prod_{b,y}(N^{(1)}_{bky})_{r_{by}}}{(m_1)_{|r|}}
\right]
=
\prod_{b,y}q_{by}^{r_{by}} .
\]
Because the arm sums are finite sums of integrable factorial monomials, termwise integration yields, for each arm \(a\),
\[
\mathbb E_{P^{\otimes\infty}}[\widehat A_{a,k}]
=
\sum_{j=0}^{M-2}\sum_{t=0}^{j}\sum_{c\in\{0,1\}^2}
B^{-1}g_{M,j}(B^{-1})^j\binom jt\,
q^{\,r^{a,c,j,t}} .
\]
% lean: integral_factorialMonomial_trunc

For fixed \(a\) and \(j\), the inner finite sum collapses as
\[
\sum_{t=0}^{j}\sum_{c\in\{0,1\}^2}
\binom jt\,q^{\,r^{a,c,j,t}}
=
\left(\sum_{b,y}q_{by}\right)q_{a1}(q_{a0}+q_{a1})^j.
\]
Indeed, the merged-exponent monomial identity writes each summand as
\(q_c q_{a1}q_{a0}^tq_{a1}^{j-t}\); summing over \(c\) gives the total mass \(\sum_{b,y}q_{by}\), and summing over \(t\) gives the binomial theorem. % lean: factorialExpansionIndex_binomial_sum

Writing \(s(q)=\sum_{b,y}q_{by}\), \(s_a(q)=q_{a0}+q_{a1}\), and \(t_a(q)=q_{a1}\), we obtain
\[
\mathbb E_{P^{\otimes\infty}}[\widehat A_{a,k}]
=
B^{-1}s(q)t_a(q)
\sum_{j=0}^{M-2}
g_{M,j}\{B^{-1}s_a(q)\}^{j}.
\]
Substituting this identity for \(a=1\) and \(a=0\) into the sparse arm representation of \(\widehat P_k\) gives
\[
\mathbb E_{P^{\otimes\infty}}[\widehat P_k]
=
B^{-1}s(q)t_1(q)\sum_{j=0}^{M-2}g_{M,j}\{B^{-1}s_1(q)\}^j
-
B^{-1}s(q)t_0(q)\sum_{j=0}^{M-2}g_{M,j}\{B^{-1}s_0(q)\}^j.
\]
By the displayed definition of \(P_{M,B}\), the right-hand side is
\(P_{M(n),B(n)}(q_k)\). % lean: integral_factorialPolynomialContribution_trunc
\end{proof}
